\documentclass[10pt,a4paper]{amsart}
\usepackage[margin=19mm]{geometry}
\usepackage{amsmath,amssymb,mathtools}
\usepackage[scr=rsfs,cal=euler]{mathalfa}
\usepackage{xfrac}
\usepackage[unicode,hidelinks,bookmarks=false]{hyperref}
\newcommand{\ie}{i.e.\ }
\newcommand{\cf}{cf.\ }
\newcommand{\resp}{respectively\ }
\newcommand{\Subsection}{\subsection}
\providecommand{\nobreakdash}{\nobreak\hbox{-}}
\setlength{\emergencystretch}{2em}
\multlinegap0pt
\usepackage{bm}
\usepackage{bbm}
\usepackage{dsfont}
\usepackage{xspace}
\usepackage{cancel}
\usepackage{mathtools}
\usepackage[shortlabels]{enumitem}
\usepackage{amsthm}
\makeatletter
\@ifundefined{theorem}{\newtheorem{theorem}{Theorem}}{}
\@ifundefined{definition}{\newtheorem{definition}[theorem]{Definition}}{}
\@ifundefined{lemma}{\newtheorem{lemma}[theorem]{Lemma}}{}
\@ifundefined{obs}{\newtheorem{obs}[theorem]{Observation}}{}
\makeatother
\let\tilde\widetilde

\newcommand{\buf}{\mathrm{buf}}
\newcommand{\eps}{\varepsilon}
\newcommand{\fail}{\mathrm{fail}}
\newcommand{\itmarab}[1]{\mbox{\rm 
({\it #1}\,\arabic{*}\hspace{0.05em})}}
\title[Robustness of the Sauer-Spencer Theorem]{Robustness of the Sauer-Spencer Theorem}
\author{Peter Allen, Julia B\"ottcher, Yoshiharu Kohayakawa, Mihir Neve}
\date{Source: arXiv:2507.03676v1 (CC BY 4.0)}
\begin{document}
\maketitle

\noindent\emph{Source and licence.} Robustness of the Sauer-Spencer Theorem. Version arXiv:2507.03676v1 (CC BY 4.0), distributed under the Creative Commons Attribution 4.0 International licence, as recorded in the arXiv licence metadata for this exact version and shown on the abstract page https://arxiv.org/abs/2507.03676v1. Full rights notice: licenses/arxiv-2507.03676-CC-BY-4.0.txt.\par\smallskip

\noindent\textit{Editorial scope.} Counted unit: the Lemma whose source label is \texttt{lem:buf\_spread\_measure} in the article (source0.tex lines 879--882, source label \texttt{lem:buf\_spread\_measure}), delivered together with the article's own complete original proof of it (source0.tex lines 883--929, one \texttt{proof} environment of 47 lines). No mathematics is written, repaired or re-derived: statement and proof are the article's own text line for line. Required context retained uncounted: thm:spread\_dense\_blowup (lines 663--678); itm:FB1 (lines 851--855); itm:FB2 (lines 851--855); itm:FB3 (lines 851--855); defn:coupled\_measure (lines 861--867); obs:coupling\_prob (lines 868--875). Every definition, display and label of those lines is retained unchanged. Omitted from this package and not redistributed: 1--662, 679--850, 856--860, 876--878, 930--1163. Nothing inside a retained excerpt is omitted or altered: every retained body is delivered exactly as the article's lines are written, so no omission receipt of any kind accompanies this package. Citations: the retained excerpts carry no citation command at all, so no bibliography entry of the article is transcribed and no reference is redistributed. Reference adaptation: none was needed and none was made; every reference inside the delivered unit resolves inside it, so no unresolved reference or citation is produced. Printed slips kept literal: no printed word, symbol, hypothesis or number of the counted statement or of its proof was altered. Layout and class: the article's own class is \texttt{article} with packages including geometry, amsmath, amssymb, amsthm, mathtools, bbm, dsfont, enumitem, microtype, parskip, graphicx, algorithm2e, hyperref, backref; the wrapper is the lane's pinned \texttt{amsart} editorial wrapper at 10pt on A4, which restates the theorem-like environments the retained text needs and adds the article's own macro definitions that the retained text needs (\texttt{\textbackslash buf}, \texttt{\textbackslash eps}, \texttt{\textbackslash fail}, \texttt{\textbackslash itmarab}, \texttt{\textbackslash tilde}), with extra wrapper packages (bm, bbm, dsfont, xspace, cancel, mathtools, enumitem). The delivered numbering is the unit's own. Assets: no retained span contains a figure environment, a graphics-inclusion command, a PDF-page inclusion, a table or a TikZ picture; no image file is copied into this package, the package declares no asset dependency and no image file is redistributed with it. Licence: version arXiv:2507.03676v1 of this article is distributed under the Creative Commons Attribution 4.0 International licence, as recorded in the arXiv licence metadata for this exact version and shown on the abstract page https://arxiv.org/abs/2507.03676v1. A full rights notice is delivered at licenses/arxiv-2507.03676-CC-BY-4.0.txt. Characters: every retained excerpt is pure ASCII, so the delivered engine, XeLaTeX with its default Latin Modern text font, reports no missing character.
\begin{theorem}[Spread Blow-up Lemma] \label{thm:spread_dense_blowup}
  For all $\Delta\ge 2$, $\Delta_{R'} \in \mathbb{N}$, $\alpha, \zeta, d>0$, and ${\kappa>1}$, there exists $\eps, \rho > 0$ such that for all $r_1 > 0$, there exists a constant $C_S > 0$ such that the following holds: 
  
  Let $R$ be a graph on $r\le r_1$ vertices and let $R'\subseteq R$ be a spanning subgraph with $\Delta(R')\leq \Delta_{R'}$. Let $H$ and $G$ be $n$-vertex graphs, given with $\kappa$-balanced, size-compatible vertex partitions $\mathcal{X}=\{X_i\}_{i\in[r]}$ and $\mathcal{V}=\{V_i\}_{i\in[r]}$, respectively, which have parts of size at
  least $n/(\kappa r_1)$. Let 
  $\tilde{\mathcal{X}}=\{\tilde{X}_i\}_{i\in[r]}$ be a family of subsets of $V(H)$, and $\mathcal{I} = \{I_x\}_{x \in V(H)}$ be a family of subsets of $V(G)$.  
  Suppose that
  \begin{enumerate}[label=\itmarab{BL}]
  \item \label{itm:SpBL1} the partition $(G,\mathcal{V})$ is $(\eps,d)$-regular on $R$ and $(\eps,d)$-super-regular on $R'$,
  \item \label{itm:SpBL2} $\Delta(H)\leq \Delta$ and $(H,\mathcal{X})$ is an $R$-partition,
  \item \label{itm:SpBL3} $\tilde{\mathcal{X}}$ is an $(\alpha,R')$-buffer for $(H,\mathcal{X})$, and
  \item \label{itm:SpBL4} $\mathcal{I} = \{I_x\}_{x \in V(H)}$ is a family of $(\rho, \zeta)$-image-restrictions. 
  \end{enumerate}

  Then there exists a $(C_S/n)$-vertex-spread probability measure on the set of all embeddings ${\varphi: V(H) \to V(G)}$ of $H$ into $G$, for which $\varphi(x) \in I_x$ for all $x \in V(H)$.
\end{theorem}

\begin{enumerate}[label=\itmarab{FB}]
    \item \label{itm:FB1} For each $x \in A_i$, we have $d_{F_i}(x) \geq  \frac{1}{2} d^b |B_i|$,
    \item \label{itm:FB2} for each $v \in B_i$, we have $d_{F_i}(v) \geq  (d^{\Delta}/100)^b |A_i|$, and
    \item \label{itm:FB3}for every set $W \subseteq B_i$ of size at least $\frac{\rho}{\mu}|B_i|$, there are at most $\frac{\rho}{\mu}|A_i|$ vertices in $A_i$ with fewer than $\frac{1}{2} d^b |W|$ candidates in $W$.
\end{enumerate}

\begin{definition}
    \label{defn:coupled_measure}
    Let $F_i = (A_i \sqcup B_i, E(F_i))$ be a bipartite graph and let~$C$ be a large constant. Let~$Z_1$ be the \emph{binomial random subgraph} of~$F_i$ generated by retaining each edge of~$E(F_i)$ independently with probability~$C/\lambda$, where $\lambda = |V(F_i)|$. Similarly, let~$Z_2$ be the \emph{uniform random subgraph} of~$F_i$ obtained by choosing~$C$ neighbours uniformly at random with replacement for each~$v \in V(F_i)$. Let~$\mathds{P}_{Z_1}$ and~$\mathds{P}_{Z_2}$ denote the probability measures on~$Z_1$ and~$Z_2$ respectively. 
    
    Then, we define $Z = Z_1 \cup Z_2$ to be a random subgraph of~$F_i$ with the probability measure 
    $$ \mathds{P}_Z(Z) \coloneqq \sum_{(Z_1, Z_2)} \mathds{1}_{Z = Z_1 \cup Z_2}(Z_1, Z_2) \cdot  \mathds{P}_{Z_1}(Z_1)\, \mathds{P}_{Z_2}(Z_2),$$ where the \emph{indicator function} $\mathds{1}_A(\mathbf{x})$ equals $1$ if $A(\mathbf{x})$ is true, and $0$ otherwise.
\end{definition} 

\begin{obs}
    \label{obs:coupling_prob}
    It can be easily verified that $\mathds{P}_Z$ is a probability measure. Further, if $\mathcal{E}$ is a monotone decreasing property on graphs, then $\mathds{P}_Z(Z \in \mathcal{E}) \leq \min\big\{ \mathds{P}_{Z_1}(Z_1 \in \mathcal{E}), \mathds{P}_{Z_2}(Z_2 \in \mathcal{E})\big\}$. Indeed, if $H \subseteq G$, then $\mathds{1}_{\mathcal{E}}(G) \leq \mathds{1}_{\mathcal{E}}(H)$ as~$\mathcal{E}$ is decreasing, and hence for $i \in \{1, 2\}$, we have
    \begin{eqnarray*}
        \mathds{P}_Z(Z \in \mathcal{E}) &=& \sum_{(Z_1, Z_2)} \mathds{1}_{\mathcal{E}}(Z_1 \cup Z_2) \cdot  \mathds{P}_{Z_1}(Z_1) \mathds{P}_{Z_2}(Z_2) \leq \sum_{(Z_1, Z_2)} \mathds{1}_{\mathcal{E}}(Z_i) \cdot \mathds{P}_{Z_1}(Z_1) \mathds{P}_{Z_2}(Z_2)\\
        &\leq& \sum_{Z_i} \Bigl(\mathds{1}_{\mathcal{E}}(Z_i) \cdot \mathds{P}_{Z_i}(Z_i)\cdot \sum_{Z_{3-i}} \mathds{P}_{Z_{3-i}}(Z_{3-i})\Bigr) = \mathds{P}_{Z_i}(Z_i \in \mathcal{E}). 
    \end{eqnarray*}   
\end{obs}

\begin{lemma}
    \label{lem:buf_spread_measure}
    Let $F$ be the auxiliary bipartite graph described above, which is a vertex-disjoint union of bipartite graphs~$\{F_i\}_{i \in [r]}$ having vertex sets~$A_i \sqcup B_i = X^{\buf}_i \sqcup V^{\buf}_i$ and satisfying properties \ref{itm:FB1}--\ref{itm:FB3}. Then, there is a constant~$C_B$ independent of $n$, where $n = |V(G)| = |V(H)|$, such that there exists a $(C_B/n)$-spread measure on the set of perfect matchings of $F$.    
\end{lemma}

\begin{proof}
    It suffices to prove the existence of a $(C_B/ n)$-spread measure on the set of perfect matchings of~$F_i$ for each $i \in [r]$. Indeed, suppose such a spread measure~$\mathds{P}_i$ exists for each~$F_i$. Let $S \subseteq E(F)$ be given with $S_i = S \cap E(F_i)$. Similarly, let~$M$ be a random perfect matching in~$F$ with $M_i = M \cap E(F_i)$. Then, as $F_i$ are vertex-disjoint bipartite graphs, the product measure $\mathds{P}_M \coloneqq \prod_{i \in [r]} \mathds{P}_i$ gives the required $(C_B/n)$-spread measure on the set of perfect matchings of $F$ as 
    \[\mathds{P}_M\bigl(\{M : S \subseteq M\}\bigr) = \prod_{i \in [r]} \mathds{P}_i\bigl(\{M_i : S_i \subseteq M_i\}\bigr) \leq \prod_{i \in [r]} \bigl( C_B/n \bigr)^{|S_i|} = \bigl( C_B/n\bigr)^{|S|}.\]
    \indent Thus, we restrict ourselves to the graph~$F_i$ for some $i \in [r]$
    having parts~$A_i$, $B_i$ of size~$\lambda$. By the assumptions on the size of~$V_i$ in the statement of Theorem~\ref{thm:spread_dense_blowup}, we have that $\lambda = |X^{\buf}_i| = \mu|X_i| \geq (\mu/\kappa r_1)n$. Let~$C$ be a large constant depending on $d$, $b$, $\kappa$, $\mu$, $\rho$, and $\Delta$, and let $(Z_1, \mathds{P}_{Z_1})$, $(Z_2, \mathds{P}_{Z_2})$, and $(Z, \mathds{P}_Z)$ be the random graphs and probability measures as in Definition~\ref{defn:coupled_measure}. 
    
    We show that the random subgraph~$Z$ satisfies Hall's theorem with probability at least~$1/2$. For this, given a positive integer~$k \in [\lambda]$, let $S \subseteq A_i$ and $T \subseteq B_i$ be vertex sets of size $|S| = k$ and $|T| = k-1$, and define~$\mathcal{A}_k$ to be the event that for some such~$S$ and~$T$, we have $N_Z(S) \subseteq T$ (that is, $S$ and $T$ witness that Hall's condition is violated). We begin by calculating the probabilities $\mathds{P}_Z(N_Z(S) \subseteq T)$ and then use them to bound the probability~$\mathds{P}_Z(\mathcal{A}_k)$ of the event $\mathcal{A}_k$ from above. 

    \noindent \textbf{Case I: $k \in \mathcal{I} \coloneqq \bigl( (2\rho/\mu)\lambda, ( 1 - \rho/\mu)\lambda\bigr)$.} For this case we rely solely on~$Z_1$. Consider the complement set $T^c = B_i\setminus T$ which has size $|T^c| > (\rho/\mu)\lambda$. Then, by~\ref{itm:FB3}, all but at most $(\rho/\mu)\lambda$ vertices of~$S$ have at least~$(d^b/2)|T^c|$ neighbours in~$T^c$ in the graph~$F_i$. Let this subset of vertices of~$S$ be denoted by~$S'$, where $|S'| \geq k - (\rho/\mu)\lambda \geq k/2$. 
    
    Hence, the number of edges in~$F_i$ between~$S'$ and~$T^c$ is at least $(k/2)\cdot(d^b/{2})|T^c| \geq (d^b\rho/4\mu) \lambda k$. Then, shifting to the binomial random graph~$Z_1$, which has edge probability $C/\lambda$, and using Observation~\ref{obs:coupling_prob} for the event $\mathcal{E} \coloneqq \{Z: N_Z(S) \subseteq T\}$, we have 
    \begin{eqnarray*}
        \mathds{P}_Z\bigl(N_Z(S) \subseteq T\bigr) &\leq& \mathds{P}_{Z_1}\!\bigl(N_{Z_1}\!(S) \subseteq T\bigr) \leq \mathds{P}_{Z_1}\!\bigl(N_{Z_1}\!(S') \subseteq T\bigr) = \mathds{P}_{Z_1}\!\bigl(E_{Z_1}\!(S', T^c) = \emptyset \bigr)\\
        &\leq& \bigl(1 - C/\lambda\bigr)^{|E_{F_i}(S', T^c)|} \leq \bigl(1 - C/\lambda\bigr)^{(d^b \rho/{4 \mu}) \lambda k} \leq e^{-C_{1}k},
    \end{eqnarray*}
    where $C_1 = (d^b \rho/4 \mu)C$. We let $C$ be large enough so that $C_1$ satisfies $e^{2-C_1}(\mu/ \rho)^2 < 1/2$. Now, by a union-bound argument and using $\binom{\lambda}{k} \leq (e\lambda/k)^k$ for large $\lambda$, we have    
    \begin{eqnarray*}
        \sum_{k \in \mathcal{I}} \mathds{P}_Z(\mathcal{A}_k) &\leq& \sum_{k \in \mathcal{I}} \binom{\lambda}{k}\binom{\lambda}{k-1} e^{-C_1k} \leq \smashoperator[r]{\sum_{k > (2\rho/\mu)\lambda}}\;\; \bigl({2e\lambda}/{k}\bigr)^{2k} e^{-C_1k}\\
        &\leq& \smashoperator{\sum_{k > (2\rho/\mu)\lambda}}\;\; \bigl(\mu^2 e^{2-C_1}/\rho^2\,\bigr)^{k} \;\leq \smashoperator[r]{\sum_{k \geq (2\rho/\mu)\lambda}}\;\; 2^{-k} \;\leq 2^{-(2\rho/\mu)\lambda} < \frac{1}{4}.
    \end{eqnarray*}
    
    \noindent \textbf{Case II: $k \leq (2\rho/\mu) \lambda$.} For this case, we rely solely on~$Z_2$. By~\ref{itm:FB1}, any vertex~$x \in S$ has degree~$d_{F_i}(x) \geq (d^b/2)\lambda$. As $|T| < k$, the probability that the event $N_{Z_2}(x) \subseteq T$ holds in~$Z_2$ is at most $\bigl(2k/d^b\lambda\bigr)^C$, which is much smaller than $1$ as $10^6\rho < d^b \mu$. Hence, for $1 \leq k \leq (2\rho/\mu)\lambda$, we have 
    \begin{eqnarray*}
        \mathds{P}_{Z}(\mathcal{A}_k) &\leq& \binom{\lambda}{k}\binom{\lambda}{k-1} \mathds{P}_Z\bigl(N_Z(S) \subseteq T\bigr) \leq \bigl(2e\lambda/{k}\bigr)^{\!2k} \,\mathds{P}_{Z_2}\bigl(N_{Z_2}(S) \subseteq T\bigr) \\
        &\leq& \biggl(\frac{2e\lambda}{k}\biggr)^{\!2k} \biggl(\frac{2k}{d^b\lambda}\biggr)^{\!Ck} = {C_2}^k\cdot \bigl(k/\lambda\bigr)^{\!(C-4)k},
    \end{eqnarray*}
    where $C_2 = 4e^2 (2/d^b)^C$. Let $C_* = C_2\cdot (2\rho/\mu)^C$, and let $C$ be large enough so that $C_* < 1$. This is possible as $4\rho < \mu d^b$. Then, letting $\mathcal{I}'$ be the interval $(\lambda^{1/3}, (2\rho/\mu)\lambda)$ and taking a union-bound over all values of $k \leq (2\rho/\mu)\lambda$, we obtain the following. 
    \begin{eqnarray*}
        \smashoperator{\sum_{k \leq (2\rho/\mu)\lambda}}\; \mathds{P}_Z(\mathcal{A}_k) &=& \sum_{k \leq \lambda^{1/3}} {C_2}^k \bigl(k/\lambda\bigr)^{k(C-4)} + \sum_{k \in \mathcal{I}'} {C_2}^k \bigl( k/\lambda\bigr)^{k(C-4)}\\
        &\leq& \sum_{k \leq \lambda^{1/3}} {\Bigl(C_2\, \lambda^{-2(C-4)/3} \Bigr)}^k + \sum_{k \in \mathcal{I}'} {\Bigl(C_2\, \bigl(2\rho/\mu\bigr)^C \Bigr)}^k \\
        &\leq& \sum_{k \leq \lambda^{1/3}} C_2\, \lambda^{-2/3} + \sum_{k \geq \lambda^{1/3}} {C_*}^k\; \leq\;  C_2\, \lambda^{-1/3} + (1- C_*)^{-1}{C_*}^{(\lambda^{1/3})} =  o_\lambda(1)
    \end{eqnarray*}
    Thus, for large values of $\lambda$, we have $\sum_{k \leq (2\rho/\mu)\lambda} \mathds{P}_Z(\mathcal{A}_k) < 1/8$. 

    \textbf{Case III: $k \geq  \big(1 - \rho/\mu\big)\lambda$.} In the random subgraph~$Z$, the event~$N_Z(S) \subseteq T$ holds if and only if we have $N_Z(T^c) \subseteq S^c$. Further, we also have $|T^c| \leq (2\rho/\mu)\lambda$ and $|S^c| = |T^c| - 1$. This resembles Case II, and so using~\ref{itm:FB2} and following the calculations as above, \emph{mutatis mutandis}, we have $\sum_{k \geq (1 - \rho/\mu)\lambda} \mathds{P}_Z(\mathcal{A}_k) < 1/8$ for large values of $\lambda$. 

    Hence, we have shown that the random graph~$Z$ fails Hall's theorem with probability at most $\sum_{k \in [\lambda]} \mathds{P}_Z(\mathcal{A}_k) \leq 1/2$, and thus with probability at least~$1/2$, the random subgraph~$Z$ contains a perfect matching. For each instance of~$Z$ which contains a perfect matching, arbitrarily fix one. This defines a random variable taking values in the set of perfect matchings of~$F_i$ together with an element~$\{\fail\}$, where $\{\fail\}$ is chosen whenever~$Z$ does not satisfy Hall's theorem. Let  $\mathds{P}_i^*$ denote the distribution of this random variable, for which the probability~$\mathds{P}_i^*\bigl(\{\fail\}\bigr)$ was shown to be at most~$1/2$. We will now show that by conditioning on not choosing the element~$\{\fail\}$, we obtain the required spread-distribution~$\mathds{P}_i$ on the set of perfect matchings of~$F_i$. 
    
    For this, let $S \subseteq E(F_i)$ be given. We can assume that~$S$ is a (not necessarily perfect) matching, for otherwise, $\mathds{P}_i\bigl(\{M : S \subseteq M\}\bigr) = 0$ and so, the spread condition is trivially satisfied. Thus, let $S$ be a matching. Observe that $\mathds{P}_i^*\bigl(\{M:S\subseteq M\}\bigr)\le\mathds{P}_Z\bigl(\{Z:S\subseteq Z\}\bigr)$ as~$M$ was chosen in~$Z$. We will bound the latter probability from above. Since~$S$ is a matching,~$Z_1$ is a binomial random subgraph, and since we selected edges of $Z_2$ independently with replacement, the events that the various edges of~$S$ are in~$Z$ are mutually independent. So let~$e \in S$ be an edge in the matching. 
    
    This edge~$e$ lies in~$Z_1$ with probability~$C/\lambda$. Similarly, the edge~$e$ lies in~$Z_2$ if either end point choses the other end point, which has probability at most $2C/d^b\lambda + (100/d^{\Delta})^b C/\lambda$. Thus for $C' = C(200/d^{\Delta})^b$, we have
    \[\mathds{P}_Z\bigl(\{Z:e\in E(Z)\}\bigr)\le \mathds{P}_{Z_1}\bigl(\{Z_1:e\in E(Z_1)\}\bigr)+\mathds{P}_{Z_2}\bigl(\{Z_2:e\in E(Z_2)\}\bigr)\le 2C'/\lambda\,.\]
    
    By independence, we have $\mathds{P}_Z\bigl(\{Z:S\subseteq Z\}\bigr)\le (2C'/\lambda)^{|S|}$, and hence this is also an upper bound on $\mathds{P}_i^*\bigl(\{M:S\subseteq M\}\bigr)$. Since the probability~$P_i^*\bigl(\{\mathrm{fail}\}\bigr)$ is at most~$1/2$, in the conditioned distribution~$\mathds{P}_i$ we have $\mathds{P}_i\bigl(\{M:S\subseteq M\}\bigr)\le 2(2C'/\lambda)^{|S|}\le (4C'/\lambda)^{|S|}$. Finally, as $\lambda = \mu |V_i| \geq \mu n/ \kappa r_1$, setting $C_B = 4\kappa r_1 C'/\mu$, it follows that $\mathds{P}_i\big(\{M: S \subseteq M\}\big) \leq (C_B/n)^{|S|}$. 
    
    Thus, for all $i \in [r]$, the probability measure~$\mathds{P}_i$ is a $(C_B/n)$-spread on the set of perfect matchings of~$F_i$, and consequently, the product measure $\mathds{P}_M \coloneqq \prod_{i \in [r]} \mathds{P}_i$ forms the required $(C_B/n)$-spread probability measure on the set of perfect matchings of~$F$.  
\end{proof}

\end{document}
